% !TeX root = ../nc-del-pezzo.tex
\section{次数1の候補の一般階数と滑らかな零方向}
\label{sec:ku-d1-candidate-geometry}

本節は次数1の多項式候補
\[
                    T_s:=\mathcal H_1+s\mathcal R_1,
                    \qquad s\in\C
\]
の幾何を調べる。
実際の $\delta_2$ との同定には，命題\ref{prop:ku-d1-polynomial-classification}の
多項式延長仮定が引き続き必要である。
候補の一般階数は $s$ によらず最大である一方，
$s\ne-7/8$ の全ての値で滑らかな曲面上の零方向を構成できる。
例外値 $s=-7/8$ については，本節では零方向の存在も不存在も結論しない。

\subsection{補正係数によらない一般階数}

\begin{proposition}[次数1の候補の一般階数3]
\label{prop:ku-d1-candidate-generic-rank}
各 $s\in\C$ に対して，滑らかな次数1曲面の非空稠密開集合上で
\[
               T_{s,X}:\Sym^2H^0(X,-K_X)\longrightarrow H^1(X,T_X)
\]
の階数は3である。
さらに一つの滑らかな曲面
\begin{equation}\label{eq:ku-d1-candidate-rank-witness}
 X_*:\quad \frac{z^2}{2}+\frac{y^3}{3}
                     +x^6+x^3u^3+u^6=0
\end{equation}
が，全ての $s$ に対する階数3の証人になる。
\end{proposition}
\begin{proof}
$B=x^6+x^3u^3+u^6$ は平方因子を持たない。
実際，$u=1$ で $x^6+x^3+1$ と
$3x^2(2x^3+1)$ に共通零点はなく，$u=0$ でも零にならない。
式\eqref{eq:ku-d1-candidate-rank-witness}の錐の臨界点では
$z=y=0$ と $B_x=B_u=0$ が必要なので，原点以外に臨界点はない。
重み付き射影空間の二つの特異点も避けるため，$X_*$ は滑らかである。

$A=0$ では $\mathcal R_1=0$ であり，
式\eqref{eq:ku-d1-reduced-jacobian}のJacobian関係は純六次式の直和因子に入る。
従って $y\Sym^4\C^2{}^*$ はそのまま商に埋め込まれ，
$T_s$ はこの因子に値を持つ $[yD^2B]/4$ である。
三つの偏微分
\[
 B_{xx}=30x^4+6xu^3,\qquad B_{xu}=9x^2u^2,
 \qquad B_{uu}=6x^3u+30u^4
\]
は一次独立である。
$x^4,x^2u^2,u^4$ の係数小行列式は $30\cdot9\cdot30=8100$ である。
sourceは3次元なので階数は3となる。
滑らかな全係数空間は整であり，束射の最大階数条件は開条件なので，
各固定 $s$ に対して主張する稠密開を得る。
\end{proof}

\begin{corollary}[多項式延長の下での一般交わりの長さ3]
\label{cor:ku-d1-conditional-generic-intersection}
実際の $\delta_2$ が命題\ref{prop:ku-d1-polynomial-classification}の意味で
共変多項式へ延長すると仮定する。
このとき，次数1の一般の滑らかな古典曲面では $\delta_2$ の階数は3であり，
Kuとshadowを同時に固定する標識付き形式交わりの局所環は
\[
                         \C[p_0,p_1]/(p_0,p_1)^2
\]
である。従って長さは3，Hilbert関数は $(1,2)$ である。
\end{corollary}
\begin{proof}
延長仮定の下では $\delta_2=T_s$ となる一つの定数 $s$ が存在する。
前命題により一般階数は3である。
これは $\Sym^2P_X$ 上で単射であるから，非零の $p$ に対し
$\delta_2(p,p)\ne0$ である。
定理\ref{thm:ku-fixed-shadow-intersection}の二次有限決定性を適用すると，
二次イデアルは二変数の全二次式で生成される。
\end{proof}

\subsection{一つの短い族による零方向}

\begin{proposition}[nilpotentな線型ベクトル場による零方向族]
\label{prop:ku-d1-short-null-family}
$t=8s+7$ と置き，
\begin{equation}\label{eq:ku-d1-short-null-family}
 \begin{split}
 A_t&=x^3u-3tu^4,\\
 B_t&=\frac{x^6}{60}-t x^3u^3+6t(15-2t)u^6,\\
 X_t&:\quad z^2/2+y^3/3+A_ty+B_t=0
 \end{split}
\end{equation}
とする。このとき全ての $t$ に対し
\[
                       T_{(t-7)/8,X_t}(\partial_x)=0
\]
がJacobian商で成立する。
曲面 $X_t$ が滑らかであるための必要十分条件は
\begin{equation}\label{eq:ku-d1-short-family-smooth-locus}
                         t\notin\{0,8,25/3,27/4\}
\end{equation}
である。この滑らかな部分では
\[
 \mathcal R_{1,X_t}(\partial_x)
                      =[-3x^4u^2-72t xu^5]\ne0
\]
である。
\end{proposition}
\begin{proof}
線型ベクトル場 $V=x\partial_u$ に対し，直接微分すると
\begin{equation}\label{eq:ku-d1-short-null-syzygy}
 \begin{split}
 2(B_t)_{xx}&=V(A_t),\\
 (4t-30)A_t(A_t)_{xx}+(20-3t)(A_t)_x^2&=V(B_t)
 \end{split}
\end{equation}
となる。一方，$y^2=-A_t$ による簡約後の $8T_s(\partial_x)$ の代表元は
\[
 2y(B_t)_{xx}+(32s-2)A_t(A_t)_{xx}-(1+24s)(A_t)_x^2
\]
なので，$t=8s+7$ を代入すると $V(A_t)y+V(B_t)=V(f_t)$ に一致する。
従ってJacobian商で零になる。

滑らかさを確認する。$u=0$ の錐の臨界点では $z=y=0$ が必要だが，
$x\ne0$ なら $(f_t)_x=x^5/10\ne0$ なので臨界点はない。
$x=u=0$ にも原点以外の臨界点はなく，曲面はambientの特異点を避ける。
従って $u=1$ のchartを調べれば十分である。
そこでは
\[
 (f_t)_x=3x^2\left(y+x^3/30-t\right),\qquad
 (f_t)_y=y^2+x^3-3t.
\]
$x=0$ では $y^2=3t$ であり，この関係の下で
\[
                  f_t=2t(45-6t-y).
\]
従って特異点があるのは $t=0$，または
\[
            (45-6t)^2-3t=3(3t-25)(4t-27)=0
\]
のときである。

$x\ne0$ の場合には $v=x^3=3t-y^2$ と書け，残る臨界条件は
\[
 h(y):=y^2-30y+27t=0.
\]
代入後の $f_t$ を $h$ で割った剰余は
\[
                    3(-3t+2y)(3t-25)
\]
である。この一次式と $h$ のresultantは
\[
                       81t(t-8)(3t-25)^2
\]
となる。$t=0$ では共通解は $y=0,v=0$ であり $x\ne0$ に反する。
$t=8$ では $y=12,v=-120$，$t=25/3$ では $y=15,v=-200$ が
実際に特異点を与える。$x=0$ の場合と合わせて
式\eqref{eq:ku-d1-short-family-smooth-locus}を得る。

最後に，表示した $\mathcal R_1$ の代表元は直接微分から従う。
$t\ne0$ では $V\mapsto V(A_t)$ は単射である。
実際，$V=(a x+b u)\partial_x+(c x+d u)\partial_u$ とすると
\[
 V(A_t)=c x^4+(3a+d)x^3u+3b x^2u^2-12ct xu^3-12dt u^4.
\]
これが零なら $c=b=d=a=0$ である。
式\eqref{eq:ku-d1-reduced-jacobian}より，純六次式がJacobian商で零になるためには
$V(A_t)=0$ かつ $V(B_t)$ がその六次式に等しい必要がある。
表示した非零六次式はこれを満たさないため，その類は非零である。
\end{proof}

\begin{corollary}[Hessian候補を判別する短い滑らかな例]
\label{cor:ku-d1-short-hessian-null}
\[
 X:\quad \frac{z^2}{2}+\frac{y^3}{3}
 +(x^3u-21u^4)y+\frac{x^6}{60}-7x^3u^3+42u^6=0
\]
は滑らかであり，
\[
 \mathcal H_{1,X}(\partial_x)=0,
 \qquad
 \mathcal R_{1,X}(\partial_x)=-3[x^4u^2+168xu^5]\ne0
\]
である。
\end{corollary}
\begin{proof}
前命題で $t=7$，従って $s=0$ とすればよい。
\end{proof}

\subsection{一つの値を除く全ての補正係数}

\begin{theorem}[ほぼ全ての候補に対する滑らかな零方向]
\label{thm:ku-d1-candidate-smooth-null}
任意の $s\in\C\setminus\{-7/8\}$ に対して，
滑らかな次数1 del Pezzo曲面 $X$ と非零の $D\in H^0(X,-K_X)$ が存在し，
\[
                             T_{s,X}(D)=0
\]
となる。
\end{theorem}
\begin{proof}
命題\ref{prop:ku-d1-short-null-family}は
$t=8s+7\notin\{0,8,25/3,27/4\}$ を扱う。
残る三つの非零値には，別の多項式族
\begin{equation}\label{eq:ku-d1-long-null-family}
 \begin{split}
 \widetilde A_t={}&x^3u-\frac{3t}{2}x^2u^2
       +\frac{3t^2}{7}xu^3-\frac{t^2(t-3)}{14}u^4,\\
 \widetilde B_t={}&-\frac{x^5u}{20}+\frac t2x^4u^2
       -\frac{t^2}{2}x^3u^3+\frac{5t^2(t-3)}{14}x^2u^4\\
       &-\frac{t^2(t^2+6t-45)}{28}xu^5
       -\frac{t^3(5t^2-87t+315)}{490}u^6
 \end{split}
\end{equation}
を使う。$\widetilde V=x\partial_x-5u\partial_u$ に対し
\[
 \begin{split}
 2(\widetilde B_t)_{xx}&=\widetilde V(\widetilde A_t),\\
 (4t-30)\widetilde A_t(\widetilde A_t)_{xx}
 +(20-3t)(\widetilde A_t)_x^2&=\widetilde V(\widetilde B_t)
 \end{split}
\]
が全 $t$ の多項式恒等式として成立する。
従ってこちらも $T_{(t-7)/8}(\partial_x)=0$ である。
$t=7$ は既存の例\eqref{eq:ku-d1-hessian-null-example}に一致する。

$t=8,25/3,27/4$ での滑らかさには，$u=1$ の判別式
\[
 P_t(x)=400\bigl(4\widetilde A_t(x,1)^3+9\widetilde B_t(x,1)^2\bigr)
\]
を使う。これは次数10，最高次係数9である。
三つの値における係数の分母は11で可逆で，$t$ の剰余はそれぞれ $8,1,4$ となる。
以下のベクトルは $x$ の昇べき順の係数を表す。
$\mathbf F_{11}[x]$ における $P_t$ の還元 $\overline P$ と
$U\overline P+W\overline P'=1$ の証明書は
\[
 \begin{aligned}
 \overline P_8&=(-3,-2,-3,0,-3,-5,0,-5,3,-5,-2),\\
 U_8&=(-4,0,1,4,-3,-4,5,0,-3),\\
 W_8&=(0,4,-2,-1,-5,1,0,-3,2,-3),
 \end{aligned}
\]
\[
 \begin{aligned}
 \overline P_1&=(1,-3,3,4,-5,-1,5,0,-4,1,-2),\\
 U_1&=(5,5,-3,-5,5,-2,-5,1,-5),\\
 W_1&=(5,3,3,5,3,-3,-3,2,4,-5),
 \end{aligned}
\]
\[
 \begin{aligned}
 \overline P_4&=(-1,-2,4,-4,-3,0,1,3,-3,0,-2),\\
 U_4&=(2,0,3,-1,-2,1,1),\\
 W_4&=(4,3,-4,-4,-5,-5,1,1)
 \end{aligned}
\]
である。従って三つの $P_t$ は $\mathbb Q$ 上平方因子を持たない。
命題\ref{prop:ku-d1-hessian-null-direction}の証明と同様に，
これは $u=1$ 上の特異点の不存在を示す。
$u=0$ での唯一の候補は $y=z=0,x\ne0$ だが，
そこで $F_u=-x^5/20\ne0$ となる。
従って三つの曲面も滑らかであり，全ての $s\ne-7/8$ を覆う。
\end{proof}

\begin{remark}[残る値と実際の二次差]
\label{rem:ku-d1-exceptional-coefficient-scope}
本節の二つの零方向族は $s=-7/8$ で特異になる。
この事実だけから，$T_{-7/8}$ が全滑らか点で全非零方向に非零であるとは言えない。
また，その性質は実際の $\delta_2$ についても未証明であり，
既存の結果はFermatを含む稠密開上の非消滅である。
仮に多項式延長と全滑らか点・全非零方向の非消滅を両方証明できれば，
前定理により係数は $s=-7/8$ に限られる。
これは二つの追加仮定を置いた帰結であり，現時点での係数決定ではない。

系\ref{cor:ku-d1-short-hessian-null}の短い曲面で実際の内在的周期を二次まで計算すれば，
多項式延長の下では $\delta_2(\partial_x,\partial_x)=s\mathcal R_1(\partial_x)$
の係数を一回の計算で判別できる。
候補の一般階数と零方向の証明書は
\path{computations/verify_degree_one_candidate_geometry.py}で再現する。
これは非可換乗法や実際の周期を計算するコードではない。
\end{remark}
